%% Supplementary: site acceptance test (Section 3.7 and 5.7 of the main text).
%% Numbers: results/v2/site_acceptance/sat_numbers.json; tables: paper/media/tables/tab_sat_*.tex,
%% both from src/v2/site_acceptance.py (evaluate, tables --plan-check).
\section{Site acceptance test: further results}
\label{sup:sat}

\paragraph{Compact tier} Fig.~\ref{fig:sat-C} and Table~\ref{tab:sat-C} repeat
the validation of the main text for the 1,600 compact-tier runs. The compact
models differ more between sites, so the development prior is wider and adds
less. The test was as safe as at the standard tier, with false acceptance of at
most 0.4\%, and its gain over local labels alone was confined to small samples.

\paragraph{Cohorts} Table~\ref{tab:sat-cohorts} breaks down false acceptance and
power at a target of 0.85 by cohort. What governs the gain is where the
development sites score relative to the target. The median prior center is an
AUROC of 0.971 on standard-tier Camelyon17 and 0.897 on MIDOG++, where the test
gains most, but 0.869 on the canine cohort and 0.861 on MIDOG~2021, where it
gains little. When the development sites score below the target, the prior is
pessimistic and costs power: on the compact-tier canine cohort the test
accepted no good site at a target of 0.85, against 5--8\% for local labels
alone. MIDOG~2021 priors rest on only two development sites.

\paragraph{Robust variant} The $\epsilon$-contaminated prior replaces the
development prior $\pi$ by $(1-\epsilon)\pi+\epsilon\,\mathcal N(0,2^2)$ on the
logit scale, with $\epsilon=0.1$; the vague component is centered at an AUROC of
0.5. The local evidence then moves the posterior faster when it contradicts the
development sites. Table~\ref{tab:sat-robust} gives its operating
characteristics. We added the variant after the first validation showed false
acceptance above 5\% in one subgroup (standard-tier Camelyon17 at a target of
0.90 with at most 50 labels), so its validation is not independent of that
observation. In that subgroup it reduced false acceptance from 8.7\% to 5.4\%
(25 labels), from 5.7\% to 4.5\% (50) and from 3.3\% to 2.8\% (100).

\paragraph{Label planner} The planner simulates equal-variance binormal scores
at a prevalence of 0.5, with a margin of 0.02 above the target and a target
power of 0.8. We drew 150 priors per tier at random and kept those with at
least one run 0.02 above a target of 0.85 (114 standard-tier and 44
compact-tier priors; 28 of the 44 are Camelyon17), and compared the planner's predicted acceptance rate of good
sites (true AUROC at least 0.02 above a target of 0.85) with the observed rate.
At the compact tier the two agreed within 5 percentage points at every sample
size. At the standard tier the planner was optimistic by 3 to 7 points
(predicted 75\% against 68\% observed at 100 labels, and 85\% against 78\% at
200). A practitioner should therefore add a margin to the planned sample size.

\paragraph{Assumptions} The development prior treats the new site as
exchangeable with the development sites. In the validation, the prior for site
$t$ uses only the other folds' results at their own held-out sites; site $t$
enters those folds only as training data, so its test labels inform neither
the prior nor the decision, but the development results are not independent of
site $t$. Local samples are drawn without replacement from the same test
patches that define the truth, which for MIDOG~2021 and 400 labels is up to a
third of them; this flatters the agreement between estimate and truth
somewhat. The decision threshold is a posterior probability; the reported
false-acceptance rates are empirical, not guaranteed. A new site that differs in a way none
of the development sites does, for example a new stain protocol, is exactly
the case that the local labels must catch. Labeled patches were drawn at
random from the site's test patches; in practice patches come in clusters from
a few slides, which increases the variance of the local AUROC, so labels should
be spread over many slides. The test concerns patch-level AUROC, not a
clinical endpoint.

\begin{figure}[pos=tbp]
  \centering
  \includegraphics[width=\linewidth]{fig_sat_C}
  \caption{Validation of the site acceptance test at the compact tier, as in
  Fig.~\ref{M-fig:sat}.}
  \label{fig:sat-C}
\end{figure}

\begin{table}[width=\linewidth,pos=tbp]
\caption{Site acceptance test at the compact tier (1,600 runs), as in
Table~\ref{M-tab:sat}.}
\label{tab:sat-C}
\footnotesize
\setlength{\tabcolsep}{3pt}
% generated by src/v2/site_acceptance.py tables (results/v2/site_acceptance/sat_rows_*.csv)
\begin{tabular*}{\tblwidth}{@{\extracolsep{\fill}}lrrrrrrrrr@{}}
\toprule
 & & \multicolumn{2}{c}{$T=0.80$} & \multicolumn{2}{c}{$T=0.85$} & \multicolumn{2}{c}{$T=0.90$} & & \\
\cmidrule(lr){3-4}\cmidrule(lr){5-6}\cmidrule(lr){7-8}
Rule & $m$ & FA (\%) & power (\%) & FA (\%) & power (\%) & FA (\%) & power (\%) & MAE & cover. (\%) \\
\midrule
Development prior only & 0 & 1.2 & 15 & 0.5 & 8 & 0.3 & 2 & -- & -- \\
\addlinespace[2pt]
Site acceptance test & 25 & 0.2 & 22 & 0.1 & 17 & 0.2 & 2 & 0.042 & 90 \\
 & 50 & 0.3 & 38 & 0.1 & 27 & 0.2 & 14 & 0.035 & 90 \\
 & 100 & 0.4 & 48 & 0.1 & 46 & 0.1 & 38 & 0.027 & 90 \\
 & 200 & 0.4 & 57 & 0.1 & 55 & 0.1 & 63 & 0.021 & 91 \\
 & 400 & 0.3 & 67 & 0.1 & 62 & 0.0 & 74 & 0.015 & 92 \\
\addlinespace[2pt]
Local labels only & 25 & 0.5 & 7 & 0.0 & 0 & 0.0 & 0 & 0.068 & 94 \\
 & 50 & 1.0 & 32 & 0.5 & 28 & 0.1 & 8 & 0.047 & 92 \\
 & 100 & 0.9 & 45 & 0.5 & 46 & 0.1 & 43 & 0.033 & 92 \\
 & 200 & 0.6 & 56 & 0.3 & 56 & 0.1 & 62 & 0.022 & 92 \\
 & 400 & 0.4 & 67 & 0.2 & 64 & 0.0 & 74 & 0.015 & 93 \\
\bottomrule
\end{tabular*}

\end{table}

\begin{table}[width=\linewidth,pos=tbp]
\caption{Robust variant ($\epsilon=0.1$) at the standard tier (top) and the
compact tier (bottom), as in Table~\ref{M-tab:sat}. The variant was added after
the first validation and is therefore not independently validated.}
\label{tab:sat-robust}
\footnotesize
\setlength{\tabcolsep}{3pt}
% generated by src/v2/site_acceptance.py tables (results/v2/site_acceptance/sat_rows_*.csv)
\begin{tabular*}{\tblwidth}{@{\extracolsep{\fill}}lrrrrrrrrr@{}}
\toprule
 & & \multicolumn{2}{c}{$T=0.80$} & \multicolumn{2}{c}{$T=0.85$} & \multicolumn{2}{c}{$T=0.90$} & & \\
\cmidrule(lr){3-4}\cmidrule(lr){5-6}\cmidrule(lr){7-8}
Rule & $m$ & FA (\%) & power (\%) & FA (\%) & power (\%) & FA (\%) & power (\%) & MAE & cover. (\%) \\
\midrule
Robust variant & 25 & 1.4 & 58 & 0.6 & 28 & 0.2 & 14 & 0.029 & 93 \\
 & 50 & 1.7 & 72 & 0.8 & 47 & 0.3 & 27 & 0.024 & 92 \\
 & 100 & 1.2 & 83 & 0.7 & 62 & 0.2 & 43 & 0.019 & 92 \\
 & 200 & 0.8 & 89 & 0.6 & 71 & 0.2 & 51 & 0.015 & 92 \\
 & 400 & 0.4 & 93 & 0.3 & 79 & 0.1 & 60 & 0.011 & 93 \\
\bottomrule
\end{tabular*}
\\[6pt]
% generated by src/v2/site_acceptance.py tables (results/v2/site_acceptance/sat_rows_*.csv)
\begin{tabular*}{\tblwidth}{@{\extracolsep{\fill}}lrrrrrrrrr@{}}
\toprule
 & & \multicolumn{2}{c}{$T=0.80$} & \multicolumn{2}{c}{$T=0.85$} & \multicolumn{2}{c}{$T=0.90$} & & \\
\cmidrule(lr){3-4}\cmidrule(lr){5-6}\cmidrule(lr){7-8}
Rule & $m$ & FA (\%) & power (\%) & FA (\%) & power (\%) & FA (\%) & power (\%) & MAE & cover. (\%) \\
\midrule
Robust variant & 25 & 0.2 & 19 & 0.1 & 13 & 0.1 & 2 & 0.043 & 91 \\
 & 50 & 0.2 & 36 & 0.1 & 25 & 0.1 & 12 & 0.035 & 91 \\
 & 100 & 0.4 & 48 & 0.1 & 45 & 0.1 & 38 & 0.027 & 91 \\
 & 200 & 0.4 & 57 & 0.1 & 54 & 0.1 & 63 & 0.021 & 91 \\
 & 400 & 0.3 & 67 & 0.1 & 62 & 0.0 & 74 & 0.015 & 92 \\
\bottomrule
\end{tabular*}

\end{table}

\begin{table}[width=\linewidth,pos=tbp]
\caption{False acceptance (\%) / power (\%) by cohort at a target of 0.85, with
50 and 100 local labels, for the site acceptance test, its robust variant and
local labels alone. Bad: number of runs whose true site AUROC is below 0.85.}
\label{tab:sat-cohorts}
\footnotesize
\setlength{\tabcolsep}{3pt}
% generated by src/v2/site_acceptance.py tables; cells: false acceptance (%) / power (%), T = 0.85
\begin{tabular*}{\tblwidth}{@{\extracolsep{\fill}}llrrrrrrr@{}}
\toprule
Tier & Cohort & bad & Site 50 & Site 100 & Robust 50 & Robust 100 & Local 50 & Local 100 \\
\midrule
S & Camelyon17 & 14 & 2.5 / 87 & 0.9 / 95 & 1.9 / 81 & 0.7 / 93 & 0.3 / 60 & 0.2 / 88 \\
 & Canine cSCC & 99 & 1.5 / 27 & 1.2 / 45 & 1.4 / 26 & 1.1 / 44 & 0.9 / 26 & 0.7 / 45 \\
 & MIDOG 2021 & 62 & 0.0 / 5 & 0.2 / 14 & 0.0 / 4 & 0.2 / 13 & 0.4 / 4 & 0.7 / 11 \\
 & MIDOG++ & 73 & 0.7 / 50 & 0.6 / 65 & 0.6 / 46 & 0.5 / 63 & 0.3 / 15 & 0.3 / 34 \\
\addlinespace[3pt]
C & Camelyon17 & 114 & 0.8 / 49 & 0.7 / 80 & 0.7 / 46 & 0.6 / 79 & 0.4 / 49 & 0.4 / 76 \\
 & Canine cSCC & 376 & 0.0 / 0 & 0.0 / 0 & 0.0 / 0 & 0.0 / 0 & 0.4 / 5 & 0.3 / 8 \\
 & MIDOG 2021 & 173 & 0.0 / 3 & 0.0 / 10 & 0.0 / 2 & 0.0 / 9 & 0.3 / 3 & 0.3 / 9 \\
 & MIDOG++ & 387 & 0.0 / 3 & 0.0 / 9 & 0.0 / 2 & 0.0 / 9 & 0.7 / 6 & 0.7 / 15 \\
\bottomrule
\end{tabular*}

\end{table}
